% Supply bound on junction vertex b

\subsubsection{Supply bound on the junction vertex \texorpdfstring{$b$}{b}}\label{sec:b-capacity}

Fix the right tail $P_R=(b,z_1,\ldots,u_R)$, a path of length~$s_2$ from~$b$ to~$u_R$.
If $p(u_R)=0$ and some sequence from~$p$ pebbles~$u_R$, then
Lemma~\ref{lem:path-lb} applied to~$P_R$ yields a reachable configuration with
$p(b)\ge 2^{s_2}$.

\begin{lemma}[Left-tail potential]\label{lem:left-potential}
Label $P_L=(v_0,v_1,\ldots,v_{s_1})$ with $v_0=u_L$ and $v_{s_1}=a$.
For a configuration~$p$, define
\[
\Phi(p)=\sum_{i=0}^{s_1} p(v_i)\,2^i.
\]
Then $\Phi(p)$ never increases under moves on~$P_L\setminus\{a\}$ or on the
right side of~$G$, and each move $a\to c$, $a\to b$, or $a\to v_{s_1-1}$ (back into
the tail) satisfies $\Delta\Phi<0$; in particular $\Delta\Phi=-2^{s_1+1}$ for
$a\to c$ and $a\to b$.
\end{lemma}

\begin{proof}
If a move occurs on an edge of~$P_L$ not incident with~$a$, the calculation in
Lemma~\ref{lem:path-lb} (with weights $2^i$ toward~$a$) shows that $\Phi$ does
not increase.
The move $v_{s_1-1}\to a$ (toward~$a$ along the tail) leaves $\Phi$ unchanged,
matching the general pattern of Lemma~\ref{lem:path-lb}.
Each move $a\to c$ or $a\to b$ removes two pebbles from~$a$ and therefore
decreases $\Phi$ by $2\cdot 2^{s_1}=2^{s_1+1}$.
The move $a\to v_{s_1-1}$ (back into the tail, when $s_1\ge 2$) decreases $\Phi$ by
$2^{s_1+1}-2^{s_1-1}>0$.
Together with the standard pattern on $P_L\setminus\{a\}$ (Lemma~\ref{lem:path-lb}) and
the move $v_{s_1-1}\to a$, the cases above cover every pebbling move incident with~$a$.
Moves on the right side of~$G$ do not change~$\Phi$.
\end{proof}

\begin{lemma}[$b$-potential bound]\label{lem:b-potential}
Let $K=2^{s_1+1}$ and $N=2^d-1$.
Define $\Xi(p)=\Phi(p)+Kp(b)$.
Every configuration reachable from~$D^*$ satisfies $\Xi(p)\le N$.
\end{lemma}

\begin{proof}
Since $\Xi(D^*)=N$, it suffices to show that no single pebbling move raises $\Xi$
above~$N$.
Let $p^-$ be the configuration immediately before a move~$m$ that yields~$p$,
with $\Xi(p^-)\le N$ by induction on sequence length.

A move $a\to b$ lowers $\Phi$ by~$K$ and raises $p(b)$ by~$1$, so $\Xi$ is unchanged.
Each move $a\to c$ or $b\to z_1$ lowers $\Xi$ by at least~$K$ (in fact by~$2K$
for each $b\to z_1$), and each move $b\to c$ lowers $\Xi$ by~$2K$.
A move $b\to a$ lowers $\Phi$ by $2^{s_1}=K/2$ and $p(b)$ by~$2$, so
$\Delta\Xi=-\tfrac{3}{2}K<0$.
All other moves except those listed below leave $\Xi$ unchanged or decrease it.

\smallskip\noindent
\textbf{Move $c\to b$.}
This requires $p^-(c)\ge 2$.
Since $D^*(c)=1$, the value $p(c)=2$ can arise only after a \emph{supply move}
$a\to c$ or $b\to c$ that raised $p(c)$ from~$0$ to~$1$ or from~$1$ to~$2$.
Each supply move lowers $\Xi$ by at least~$K$ ($a\to c$) or~$2K$ ($b\to c$).
Apply the inductive hypothesis to the state immediately before the \emph{last}
supply move that enabled $p^-(c)\ge 2$: there $\Xi\le N-K$.
The move $c\to b$ lowers $\Phi$ by~$K$ and raises $p(b)$ by~$1$, so
$\Xi(p)=\Xi(p^-)-K+K=\Xi(p^-)\le N-K<N$.
If $p(c)\ge 2$ because of several supply cycles, the same bound applies to the
prefix ending at the supply move that immediately preceded~$m$.

\smallskip\noindent
\textbf{Move $c\to a$.}
This raises $\Xi$ by $K/2=2^{s_1}$.
As above, $p(c)\ge 2$ requires a preceding supply move $a\to c$ or $b\to c$
with $\Xi\le N-K$ at the state just before that supply move, so
$\Xi(p)\le N-K+K/2<N$.

\smallskip\noindent
\textbf{Move $z_1\to b$.}
This requires $p^-(z_1)\ge 2$.
Each pebble on~$z_1$ arrived from~$b$ via a move $b\to z_1$, and each such move
lowers $\Xi$ by~$2K$.
To accumulate $p(z_1)\ge 2$ requires at least two moves $b\to z_1$, lowering
$\Xi$ by at least~$4K$ before $z_1\to b$ can occur; the subsequent $z_1\to b$
raises $\Xi$ by only~$K$, so $\Xi(p)\le N$ is preserved.

\smallskip\noindent
In every case $\Xi(p)\le N$, completing the induction.
\end{proof}

\begin{lemma}[$b$-capacity bound]\label{lem:b-capacity}
Starting from the critical configuration~$D^*$, every reachable
configuration~$p$ satisfies
\[
p(b)\le 2^{s_2}-1.
\]
\end{lemma}

\begin{proof}
Let $K=2^{s_1+1}$ and $N=2^d-1$.
By Lemma~\ref{lem:b-potential}, $\Phi(p)+Kp(b)\le N$.
Since $\Phi(p)\ge 0$, we obtain
\[
p(b)\le\frac{N}{K}
=\frac{2^d-1}{2^{s_1+1}}
=2^{s_2}-\frac{1}{2^{s_1+1}}.
\]
Because $0<1/2^{s_1+1}<1$ and $p(b)$ is an integer, we have
$p(b)\le\lfloor N/K\rfloor=2^{s_2}-1$.
\end{proof}

\begin{corollary}[Bypass through $c$ is insufficient]\label{cor:bypass}
No sequence of pebbling moves starting from~$D^*$ can pebble~$u_R$.
\end{corollary}

\begin{proof}
By Lemma~\ref{lem:path-lb} on the right tail~$P_R$, pebbling $u_R$ from~$D^*$
would require some reachable state with $p(b)\ge 2^{s_2}$.
Lemma~\ref{lem:b-capacity} shows $p(b)\le 2^{s_2}-1$ for every reachable state.
This is a contradiction.
\end{proof}
